%% toc: P2 11 11.7
\section{Proofs for Post's five maximal classes and the functional-completeness criterion}
\label{sec:annex-p02-w02-p02-008-post}

Notation is that of Section~\ref{sec:p02-w02-p02-008-post}: truth values are
\(0,1\); \(\oplus\) is addition modulo \(2\); \(\mathcal O_n\),
\(\mathcal O\), \(e^n_i\), \(\kappa^n_b\), composition, clones and
\(\langle C\rangle\) are as in \cref{def:p02-w02-p02-008-operations}; and
\(C^+\) is the operation set of a connective set
(\cref{def:p02-w02-p02-008-arity}).  Two arithmetic facts about \(\mathbb B\)
are used throughout: \(x\oplus x=0\) and \(xx=x\) for every
\(x\in\mathbb B\); and two members of \(\mathbb B\) that are distinct are
\(y\) and \(1-y\) for some \(y\).

\emph{Formulas over a connective set.}  For a set \(C\) of truth-functional
connectives, every formula over \(C\) is a propositional variable or has the
form \(c(\varphi_1,\ldots,\varphi_k)\) for a \(k\)-ary \(c\in C\)
(\(k\geq0\); for \(k=0\) the formula is \(c\) itself).  Write \(f_c\) for the
interpreting operation of a member \(c\) of positive arity and \(b_c\) for the
value of a nullary member.  The recursive truth clauses are
\(\widehat v(p)=v(p)\), \(\widehat v(c)=b_c\) for nullary \(c\), and
\(\widehat v(c(\varphi_1,\ldots,\varphi_k))=
f_c(\widehat v(\varphi_1),\ldots,\widehat v(\varphi_k))\) for \(k\geq1\).
\begin{enumerate}
\item[(F1)] \emph{Coincidence.}  If valuations \(v,w\) agree on the variables
  of \(\varphi\), then \(\widehat v(\varphi)=\widehat w(\varphi)\).  By
  induction on \(\varphi\): for a variable this is the hypothesis; for a
  nullary \(c\) both sides are \(b_c\); for
  \(c(\varphi_1,\ldots,\varphi_k)\) the variables of each \(\varphi_j\) are
  variables of the compound, so the induction hypothesis gives
  \(\widehat v(\varphi_j)=\widehat w(\varphi_j)\) for each \(j\), and the
  truth clause gives the same value on both sides.  This is the argument of
  \cref{lem:prop-coincidence} for formulas over \(C\).  It makes the truth
  function of \(\varphi\) in \(p_1,\ldots,p_n\) well defined whenever the
  variables of \(\varphi\) are among \(p_1,\ldots,p_n\): its value at
  \(\bar a\) is \(\widehat v(\varphi)\) for any \(v\) with \(v(p_i)=a_i\)
  for \(i\leq n\).
\item[(F2)] \emph{Substitution.}  For formulas \(\varphi,\psi_1,\ldots,\psi_m\)
  over \(C\), let \(\varphi[\bar\psi/\bar p]\) be the result of replacing
  every occurrence of \(p_j\) in \(\varphi\) by \(\psi_j\), simultaneously
  for \(j=1,\ldots,m\); recursively, \(p_j\) becomes \(\psi_j\) for
  \(j\leq m\), every other variable and every nullary connective is
  unchanged, and \(c(\varphi_1,\ldots,\varphi_k)\) becomes
  \(c(\varphi_1[\bar\psi/\bar p],\ldots,\varphi_k[\bar\psi/\bar p])\).  For a
  valuation \(v\), let \(w\) agree with \(v\) except that
  \(w(p_j)=\widehat v(\psi_j)\) for \(j\leq m\).  Then
  \(\widehat v(\varphi[\bar\psi/\bar p])=\widehat w(\varphi)\).  By
  induction on \(\varphi\): for \(p_j\) with \(j\leq m\) both sides are
  \(\widehat v(\psi_j)\); for any other variable \(q\) both are
  \(v(q)=w(q)\); for nullary \(c\) both are \(b_c\); and for
  \(c(\varphi_1,\ldots,\varphi_k)\) the left side is
  \(f_c(\widehat v(\varphi_1[\bar\psi/\bar p]),\ldots)\), which by the
  induction hypothesis is \(f_c(\widehat w(\varphi_1),\ldots)=
  \widehat w(c(\varphi_1,\ldots,\varphi_k))\).  The same recursion shows that
  every variable of \(\varphi[\bar\psi/\bar p]\) is a variable of some
  \(\psi_j\) or a variable of \(\varphi\) other than \(p_1,\ldots,p_m\).
\end{enumerate}

\begin{annexitem}{annex:p02-w02-p02-008-gen}{thm:p02-w02-p02-008-generated}{Generated clones}
\begin{proof}
Let \(\mathcal K_C\) be the family of clones \(K\) with \(C\subseteq K\); it
contains \(\mathcal O\), and \(\langle C\rangle=\bigcap\mathcal K_C\).

\emph{1.}  Every projection lies in every \(K\in\mathcal K_C\), hence in
\(\langle C\rangle\).  Let \(f\in\langle C\rangle\cap\mathcal O_m\) and
\(g_1,\ldots,g_m\in\langle C\rangle\cap\mathcal O_n\).  For each
\(K\in\mathcal K_C\), \(f\) and every \(g_j\) lie in \(K\), so
\(f\circ(g_1,\ldots,g_m)\in K\) because \(K\) is a clone; hence the
composition lies in \(\bigcap\mathcal K_C\).  So \(\langle C\rangle\) is a
clone.  \(C\) is included in every member of \(\mathcal K_C\), hence in their
intersection.

\emph{2.}  If \(K\) is a clone with \(C\subseteq K\), then
\(K\in\mathcal K_C\), and an intersection is included in each of the sets
intersected.

\emph{3.}  By part~1, \(\langle D\rangle\) is a clone; if it includes \(C\),
part~2 gives \(\langle C\rangle\subseteq\langle D\rangle\).  If
\(C\subseteq D\), then \(C\subseteq D\subseteq\langle D\rangle\) by part~1.

\emph{4.}  A clone \(K\) includes \(K\), so part~2 gives
\(\langle K\rangle\subseteq K\), and part~1 gives \(K\subseteq\langle
K\rangle\).

\emph{5.}  If \(\langle C\rangle\subseteq K\), then \(C\subseteq K\) by
part~1.  The converse is part~2.
\end{proof}
\end{annexitem}

\begin{annexitem}{annex:p02-w02-p02-008-formulas}{thm:p02-w02-p02-008-formulas-clone}{Formulas and clones}
\begin{proof}
For \(n\geq1\), let \(R_n\) be the set of \(f\in\mathcal O_n\) represented by
some formula over \(C\), and \(R=\bigcup_{n\geq1}R_n\).  The theorem says
\(R=\langle C^+\rangle\).

\emph{Step 1: \(R\subseteq\langle C^+\rangle\).}  We show by induction on
\(\varphi\): for every \(n\geq1\) such that the variables of \(\varphi\) are
among \(p_1,\ldots,p_n\), the truth function \(f^n_\varphi\) of \(\varphi\)
in \(p_1,\ldots,p_n\) lies in \(\langle C^+\rangle\).  Recall that
\(\langle C^+\rangle\) is a clone including \(C^+\)
(\cref{thm:p02-w02-p02-008-generated}).
If \(\varphi\) is \(p_i\), then \(i\leq n\) and \(f^n_\varphi=e^n_i\), a
projection.  If \(\varphi\) is a nullary \(c\), then \(f^n_\varphi\) is the
constant \(\kappa^n_{b_c}\), which equals the composition
\(\kappa^1_{b_c}\circ(e^n_1)\) of the member \(\kappa^1_{b_c}\) of
\(C^+\) with a projection.  If \(\varphi\) is
\(c(\varphi_1,\ldots,\varphi_k)\) with \(k\geq1\), the variables of each
\(\varphi_j\) are among \(p_1,\ldots,p_n\), so each
\(f^n_{\varphi_j}\in\langle C^+\rangle\) by the induction hypothesis, and the
truth clause gives
\(f^n_\varphi=f_c\circ(f^n_{\varphi_1},\ldots,f^n_{\varphi_k})\) with
\(f_c\in C^+\); so \(f^n_\varphi\in\langle C^+\rangle\).

\emph{Step 2: \(R\) is a clone including \(C^+\).}  The variable \(p_i\)
represents \(e^n_i\) for every \(n\geq i\), so \(R\) contains every
projection.  If \(c\in C\) has arity \(k\geq1\), the formula
\(c(p_1,\ldots,p_k)\) has value \(f_c(a_1,\ldots,a_k)\) under any \(v\) with
\(v(p_i)=a_i\), so it represents \(f_c\).  If \(c\) is nullary, the formula
\(c\) has no variable, so its variables are among \(p_1\), and it represents
\(\kappa^1_{b_c}\).  Hence \(C^+\subseteq R\).  For composition, let
\(f\in R_m\) be represented by \(\varphi\) and \(g_1,\ldots,g_m\in R_n\) by
\(\psi_1,\ldots,\psi_m\).  Put \(\chi=\varphi[\bar\psi/\bar p]\).  The
variables of \(\varphi\) are among \(p_1,\ldots,p_m\), all of which are
replaced, so by (F2) the variables of \(\chi\) are variables of the
\(\psi_j\), hence among \(p_1,\ldots,p_n\).  Let \(\bar a\in\mathbb B^n\)
and \(v(p_i)=a_i\) for \(i\leq n\).  By (F2),
\(\widehat v(\chi)=\widehat w(\varphi)\), where
\(w(p_j)=\widehat v(\psi_j)=g_j(\bar a)\) for \(j\leq m\).  Since
\(\varphi\) represents \(f\), \(\widehat w(\varphi)=f(g_1(\bar a),\ldots,
g_m(\bar a))\).  So \(\chi\) represents \(f\circ(g_1,\ldots,g_m)\), which
therefore lies in \(R\).

\emph{Step 3.}  By step~2 and \cref{thm:p02-w02-p02-008-generated}(2),
\(\langle C^+\rangle\subseteq R\); with step~1, \(R=\langle C^+\rangle\).
Finally, \(C\) is positive-arity functionally complete iff
\(R=\mathcal O\), iff \(\langle C^+\rangle=\mathcal O\).
\end{proof}
\end{annexitem}

\begin{annexitem}{annex:p02-w02-p02-008-nullary}{thm:p02-w02-p02-008-nullary}{Nullary functions need a nullary connective}
\begin{proof}
\emph{1.}  By induction on formulas over \(C\): a variable contains a
variable; a compound \(c(\varphi_1,\ldots,\varphi_k)\) has \(k\geq1\),
because \(C\) has no nullary member, and \(\varphi_1\) contains a variable by
the induction hypothesis, so the compound does.  A formula that represents a
nullary truth function has its variables among the empty list of variables,
so it has none; by what was just shown there is no such formula.  So neither
nullary function is represented and \(C\) is not nullary-inclusive
functionally complete.

\emph{2.}  Suppose \(C\) is nullary-inclusive functionally complete.  The
condition for every \(n\geq0\) includes the condition for every \(n\geq1\),
so \(C\) is positive-arity functionally complete; and some formula
represents a nullary function, so by part~1 \(C\) has a nullary member.

Conversely, let \(C\) be positive-arity functionally complete and let
\(c\in C\) be nullary with value \(b\).  The formula \(c\) has no variable
and value \(b\) under every valuation, so it represents the nullary function
with value \(b\).  By positive-arity completeness for \(n=1\), some formula
\(\nu\) whose only possible variable is \(p_1\) represents the unary
operation \(\neg\).  Let \(\nu'=\nu[c/p_1]\) in the notation of (F2).  By
(F2) \(\nu'\) has no variable (the only possible variable \(p_1\) of
\(\nu\) is replaced, and \(c\) has none), and for every \(v\),
\(\widehat v(\nu')=\widehat w(\nu)\) where \(w(p_1)=\widehat v(c)=b\); as
\(\nu\) represents \(\neg\), this value is \(1-b\).  So \(c\) and \(\nu'\)
represent the two nullary functions, and positive arities are covered by
hypothesis.
\end{proof}
\end{annexitem}

\begin{annexitem}{annex:p02-w02-p02-008-tests}{thm:p02-w02-p02-008-tests}{Exact membership tests}
\begin{proof}
\emph{1.}  If \(f\in\mathsf S\), then \(f(\bar a^{\,c})=1-f(\bar a)\neq
f(\bar a)\) for every \(\bar a\), in particular for those with \(a_1=0\).
Conversely, assume the test and let \(\bar a\) be arbitrary.  If \(a_1=0\),
the test gives \(f(\bar a^{\,c})\neq f(\bar a)\), so
\(f(\bar a^{\,c})=1-f(\bar a)\).  If \(a_1=1\), then
\(\bar b=\bar a^{\,c}\) has \(b_1=0\) and \(\bar b^{\,c}=\bar a\); the test
for \(\bar b\) gives \(f(\bar a)\neq f(\bar a^{\,c})\), so again
\(f(\bar a^{\,c})=1-f(\bar a)\).  So \(f\in\mathsf S\).  There are
\(2^{n-1}\) tuples with \(a_1=0\).

\emph{2.}  If \(f\in\mathsf M\), the test holds because
\(\bar a\leq\bar a[i{:=}1]\).  Conversely, assume the test and let
\(\bar a\leq\bar b\).  List the coordinates where \(\bar a\) and \(\bar b\)
differ as \(i_1<\cdots<i_k\); at each of them \(a_{i_j}=0\) and
\(b_{i_j}=1\), because \(\bar a\leq\bar b\).  Put \(\bar d^{\,0}=\bar a\) and
\(\bar d^{\,j}=\bar d^{\,j-1}[i_j{:=}1]\) for \(1\leq j\leq k\).  Coordinate
\(i_j\) of \(\bar d^{\,j-1}\) is still \(a_{i_j}=0\), since only
\(i_1,\ldots,i_{j-1}\) have been changed, so the test gives
\(f(\bar d^{\,j-1})\leq f(\bar d^{\,j})\).  Also \(\bar d^{\,k}=\bar b\).
By transitivity, \(f(\bar a)\leq f(\bar b)\), so \(f\in\mathsf M\).  For the
second sentence: \(f\notin\mathsf M\) iff the test fails for some \(\bar a\)
and \(i\) with \(a_i=0\), that is, \(f(\bar a)>f(\bar a[i{:=}1])\); in
\(\mathbb B\) this means \(f(\bar a)=1\) and \(f(\bar a[i{:=}1])=0\).

\emph{3.}  If the displayed identity holds, then \(f\) is affine by
definition.  Conversely, let
\(f(\bar a)=d_0\oplus d_1a_1\oplus\cdots\oplus d_na_n\) for every \(\bar a\).
At \(\bar a=\bar0\) every product \(d_ia_i\) is \(0\), so \(d_0=f(\bar0)=c_0\).
At \(\bar a=\bar0[i{:=}1]\) only the \(i\)-th product survives, so
\(f(\bar0[i{:=}1])=d_0\oplus d_i\); adding \(d_0\) to both sides gives
\(d_i=f(\bar0[i{:=}1])\oplus f(\bar0)=c_i\).  So any coefficients that work
are \(c_0,\ldots,c_n\), and the identity with these coefficients holds.
\end{proof}
\end{annexitem}

\begin{annexitem}{annex:p02-w02-p02-008-profile}{thm:p02-w02-p02-008-profile}{Membership profile of the named operations}
\begin{proof}
\emph{Truth tables.}  For the binary operations, list the values at
\((0,0)\), \((0,1)\), \((1,0)\), \((1,1)\) in that order; each is computed from the
polynomial in Table~\ref{tab:p02-w02-p02-008-profile} using \(x\oplus x=0\):
\(\land\colon0,0,0,1\); \(\lor\colon0,1,1,1\); \(\to\colon1,1,0,1\);
\(\leftrightarrow\colon1,0,0,1\); \(\oplus\colon0,1,1,0\);
\(\mathrm{NAND}\colon1,1,1,0\); \(\mathrm{NOR}\colon1,0,0,0\).  For example
\(\to\) at \((1,1)\) is \(1\oplus1\oplus1=1\), and \(\lor\) at \((1,1)\) is
\(1\oplus1\oplus1=1\).  The unary operations take the values \(e^1_1\colon
0,1\); \(\neg\colon1,0\); \(\kappa^1_0\colon0,0\); \(\kappa^1_1\colon1,1\)
at \(0,1\).  \(\operatorname{maj}(\bar a)=1\) iff at least two coordinates of
\(\bar a\) are \(1\).

\emph{\(\mathsf T_0\) and \(\mathsf T_1\).}  Read the values at \(\bar0\) and
\(\bar1\).  At \(\bar0\): \(0\) for \(e^1_1,\kappa^1_0,\land,\lor,\oplus,
\operatorname{maj}\) and \(1\) for \(\neg,\kappa^1_1,\to,\leftrightarrow,
\mathrm{NAND},\mathrm{NOR}\).  At \(\bar1\): \(1\) for \(e^1_1,\kappa^1_1,
\land,\lor,\to,\leftrightarrow,\operatorname{maj}\) and \(0\) for
\(\neg,\kappa^1_0,\oplus,\mathrm{NAND},\mathrm{NOR}\).  These give the first
two columns.

\emph{\(\mathsf S\).}  Use \cref{thm:p02-w02-p02-008-tests}(1).  A unary \(f\)
is self-dual iff \(f(1)\neq f(0)\): true for \(e^1_1\) and \(\neg\), false
for \(\kappa^1_0\) and \(\kappa^1_1\).  A binary \(f\) is self-dual iff
\(f(1,1)\neq f(0,0)\) and \(f(1,0)\neq f(0,1)\).  This fails for \(\land\),
\(\lor\), \(\mathrm{NAND}\) and \(\mathrm{NOR}\), whose values at \((0,1)\)
and \((1,0)\) agree, and for \(\to\), \(\leftrightarrow\) and \(\oplus\),
whose values at \((0,0)\) and \((1,1)\) agree.  For \(\operatorname{maj}\):
if \(\bar a\) has \(k\) coordinates equal to \(1\), then \(\bar a^{\,c}\)
has \(3-k\); \(\operatorname{maj}(\bar a)=1\) iff \(k\geq2\), iff
\(3-k\leq1\), iff \(\operatorname{maj}(\bar a^{\,c})=0\).  So
\(\operatorname{maj}\in\mathsf S\).

\emph{\(\mathsf M\).}  Since values lie in \(\mathbb B\), \(f\) is monotone
iff \(\bar a\leq\bar b\) and \(f(\bar a)=1\) imply \(f(\bar b)=1\).  Constant
operations and \(e^1_1\) are monotone.  If \(\land(\bar a)=1\), both
coordinates of \(\bar a\) are \(1\), hence both of \(\bar b\), so
\(\land(\bar b)=1\).  If \(\lor(\bar a)=1\), some coordinate of \(\bar a\),
hence of \(\bar b\), is \(1\).  If \(\operatorname{maj}(\bar a)=1\), two
coordinates of \(\bar a\), hence of \(\bar b\), are \(1\).  The others fail
at a comparable pair: \(\neg\) at \(0\leq1\) (\(1>0\)); \(\to\),
\(\leftrightarrow\) and \(\mathrm{NOR}\) at \((0,0)\leq(1,0)\) (values
\(1\) and \(0\)); \(\oplus\) at \((1,0)\leq(1,1)\) (values \(1\) and \(0\));
\(\mathrm{NAND}\) at \((0,0)\leq(1,1)\) (values \(1\) and \(0\)).

\emph{\(\mathsf L\).}  Use \cref{thm:p02-w02-p02-008-tests}(3).  For unary
\(f\), \(c_0=f(0)\) and \(c_1=f(1)\oplus f(0)\), and the identity holds at
\(0\) (value \(c_0\)) and at \(1\) (value \(c_0\oplus c_1=f(1)\)); so every
unary operation is affine.  For binary \(f\), the identity holds at
\((0,0),(1,0),(0,1)\) by the choice of \(c_0,c_1,c_2\), and at \((1,1)\) it
reads \(f(1,1)=c_0\oplus c_1\oplus c_2=f(1,0)\oplus f(0,1)\oplus f(0,0)\);
adding \(f(1,1)\) to both sides, a binary \(f\) is affine iff the sum of its
four values is \(0\), that is, iff it takes the value \(1\) an even number
of times.  From the truth tables, \(\leftrightarrow\) and \(\oplus\) take it
twice and are affine; \(\land\) and \(\mathrm{NOR}\) once, \(\lor\), \(\to\)
and \(\mathrm{NAND}\) three times, and these are not.  For
\(\operatorname{maj}\), \(c_0=\operatorname{maj}(0,0,0)=0\) and
\(c_i=\operatorname{maj}(\bar0[i{:=}1])\oplus0=0\), so the identity would
make \(\operatorname{maj}\) constantly \(0\); but
\(\operatorname{maj}(1,1,1)=1\).  So \(\operatorname{maj}\notin\mathsf L\).
\end{proof}
\end{annexitem}

\begin{annexitem}{annex:p02-w02-p02-008-five}{thm:p02-w02-p02-008-five-clones}{Five proper clones}
\begin{proof}
\emph{Projections.}  For \(e^n_i\): \(e^n_i(\bar0)=0\), \(e^n_i(\bar1)=1\),
\(e^n_i(\bar a^{\,c})=1-a_i=1-e^n_i(\bar a)\), \(\bar a\leq\bar b\) gives
\(a_i\leq b_i\), and \(e^n_i(\bar a)=0\oplus a_i\) is affine with
\(c_i=1\) and every other coefficient \(0\).  So every projection lies in all
five classes.

\emph{Closure under composition.}  Let \(f\in\mathcal O_m\),
\(g_1,\ldots,g_m\in\mathcal O_n\), \(h=f\circ(g_1,\ldots,g_m)\), and write
\(\bar g(\bar a)=(g_1(\bar a),\ldots,g_m(\bar a))\).
\begin{itemize}
\item \(\mathsf T_0\): if all lie in \(\mathsf T_0\), then
  \(\bar g(\bar0)=\bar0\) and \(h(\bar0)=f(\bar0)=0\).  \(\mathsf T_1\) is
  the same with \(\bar1\).
\item \(\mathsf S\): if all lie in \(\mathsf S\), then
  \(\bar g(\bar a^{\,c})=(1-g_1(\bar a),\ldots,1-g_m(\bar a))=
  \bar g(\bar a)^{\,c}\), so \(h(\bar a^{\,c})=f(\bar g(\bar a)^{\,c})=
  1-f(\bar g(\bar a))=1-h(\bar a)\).
\item \(\mathsf M\): if all lie in \(\mathsf M\) and \(\bar a\leq\bar b\),
  then \(g_j(\bar a)\leq g_j(\bar b)\) for each \(j\), so
  \(\bar g(\bar a)\leq\bar g(\bar b)\) and \(h(\bar a)\leq h(\bar b)\).
\item \(\mathsf L\): if \(f(\bar y)=c_0\oplus\bigoplus_{j=1}^mc_jy_j\) and
  \(g_j(\bar a)=d_{j0}\oplus\bigoplus_{i=1}^nd_{ji}a_i\), then, by
  distributivity of multiplication over \(\oplus\) and the associativity and
  commutativity of \(\oplus\),
  \[
   h(\bar a)=c_0\oplus\bigoplus_{j=1}^mc_j\Bigl(d_{j0}\oplus
   \bigoplus_{i=1}^nd_{ji}a_i\Bigr)
   =\Bigl(c_0\oplus\bigoplus_{j=1}^mc_jd_{j0}\Bigr)\oplus
   \bigoplus_{i=1}^n\Bigl(\bigoplus_{j=1}^mc_jd_{ji}\Bigr)a_i,
  \]
  which is affine.
\end{itemize}
So each class is a clone.

\emph{Properness and non-inclusion.}  By
\cref{thm:p02-w02-p02-008-profile}, for each ordered pair of distinct
classes \(K,K'\) the following operation lies in \(K\) and not in \(K'\):
\begin{center}\small
\begin{tabular}{@{}l|ccccc@{}}
\(K\setminus K'\) & \(K'=\mathsf T_0\) & \(\mathsf T_1\) & \(\mathsf S\) & \(\mathsf M\) & \(\mathsf L\)\\\hline
\(K=\mathsf T_0\) & & \(\kappa^1_0\) & \(\kappa^1_0\) & \(\oplus\) & \(\land\)\\
\(\mathsf T_1\) & \(\kappa^1_1\) & & \(\kappa^1_1\) & \(\to\) & \(\land\)\\
\(\mathsf S\) & \(\neg\) & \(\neg\) & & \(\neg\) & \(\operatorname{maj}\)\\
\(\mathsf M\) & \(\kappa^1_1\) & \(\kappa^1_0\) & \(\kappa^1_0\) & & \(\land\)\\
\(\mathsf L\) & \(\neg\) & \(\neg\) & \(\kappa^1_0\) & \(\neg\) &
\end{tabular}
\end{center}
Each entry is read from Table~\ref{tab:p02-w02-p02-008-profile}: for example
\(\oplus\) is in \(\mathsf T_0\) and not in \(\mathsf M\), and
\(\operatorname{maj}\) is in \(\mathsf S\) and not in \(\mathsf L\).  So no
class is included in another.  Each class is proper: \(\neg\) lies outside
\(\mathsf T_0\), \(\mathsf T_1\) and \(\mathsf M\), \(\kappa^1_0\) outside
\(\mathsf S\), and \(\land\) outside \(\mathsf L\).
\end{proof}
\end{annexitem}

\begin{annexitem}{annex:p02-w02-p02-008-neg-and}{thm:p02-w02-p02-008-neg-and}{Negation and conjunction generate every operation}
\begin{proof}
Let \(K=\langle\{\neg,\land\}\rangle\), a clone
(\cref{thm:p02-w02-p02-008-generated}).  For \(u,w\in K\cap\mathcal O_n\)
the compositions \(\neg\circ(u)\) and \(\land\circ(u,w)\) lie in \(K\); write
them \(\neg u\) and \(u\land w\).  Also
\(u\lor w\defeq\neg(\neg u\land\neg w)\) lies in \(K\), and
\((u\lor w)(\bar a)=1\oplus(1\oplus u(\bar a))(1\oplus w(\bar a))\), which is
\(0\) iff \(u(\bar a)=w(\bar a)=0\); so \((u\lor w)(\bar a)=1\) iff
\(u(\bar a)=1\) or \(w(\bar a)=1\).  Likewise \((u\land w)(\bar a)=1\) iff
both are \(1\).

Let \(f\in\mathcal O_n\).  If \(f(\bar a)=0\) for every \(\bar a\), then
\(f=e^n_1\land\neg e^n_1\), whose value \(a_1(1\oplus a_1)=a_1\oplus a_1\) is
\(0\).  Otherwise list the tuples with value \(1\) as
\(\bar a^{\,1},\ldots,\bar a^{\,r}\), \(r\geq1\).  For a tuple \(\bar a\) and
\(i\leq n\), let \(\ell^{\bar a}_i=e^n_i\) if \(a_i=1\) and
\(\ell^{\bar a}_i=\neg e^n_i\) if \(a_i=0\); then
\(\ell^{\bar a}_i(\bar b)=1\) iff \(b_i=a_i\).  Put
\(m_{\bar a}=\ell^{\bar a}_1\) if \(n=1\) and
\(m_{\bar a}=(\cdots(\ell^{\bar a}_1\land\ell^{\bar a}_2)\land\cdots)\land
\ell^{\bar a}_n\) if \(n\geq2\).  By the property of \(\land\) and induction
on the number of conjuncts, \(m_{\bar a}(\bar b)=1\) iff \(b_i=a_i\) for
every \(i\), iff \(\bar b=\bar a\).  Put \(d=m_{\bar a^1}\) if \(r=1\) and
\(d=(\cdots(m_{\bar a^1}\lor m_{\bar a^2})\lor\cdots)\lor m_{\bar a^r}\) if
\(r\geq2\).  By the property of \(\lor\) and induction on the number of
disjuncts, \(d(\bar b)=1\) iff \(\bar b=\bar a^{\,j}\) for some \(j\), iff
\(f(\bar b)=1\).  So \(f=d\in K\).  Hence \(K=\mathcal O\).
\end{proof}
\end{annexitem}

\begin{annexitem}{annex:p02-w02-p02-008-anf}{thm:p02-w02-p02-008-anf}{Algebraic normal form}
\begin{proof}
Write \(\Sigma_{\mathcal J}(\bar a)=\bigoplus_{J\in\mathcal J}\prod_{i\in
J}a_i\).  Existence is by induction on \(n\).

For \(n=1\), \(f(x)=f(0)\oplus(f(0)\oplus f(1))x\): at \(x=0\) the right side
is \(f(0)\), and at \(x=1\) it is \(f(0)\oplus f(0)\oplus f(1)=f(1)\).  Take
\(\mathcal J\) to contain \(\varnothing\) iff \(f(0)=1\) and \(\{1\}\) iff
\(f(0)\oplus f(1)=1\).

For the step, let \(f\in\mathcal O_{n+1}\) and define
\(f_0,f_1\in\mathcal O_n\) by \(f_0(\bar a)=f(\bar a,0)\) and
\(f_1(\bar a)=f(\bar a,1)\).  Then
\(f(\bar a,x)=f_0(\bar a)\oplus x\bigl(f_0(\bar a)\oplus f_1(\bar a)\bigr)\),
checked at \(x=0\) and \(x=1\) as before.  By the induction hypothesis,
\(f_0=\Sigma_{\mathcal J_0}\) and \(f_0\oplus f_1=\Sigma_{\mathcal J_1}\) for
sets \(\mathcal J_0,\mathcal J_1\) of subsets of \(\{1,\ldots,n\}\).
Distributing the factor \(x\) over the sum,
\[
 f(\bar a,x)=\bigoplus_{J\in\mathcal J_0}\prod_{i\in J}a_i\ \oplus
 \bigoplus_{J\in\mathcal J_1}\Bigl(\prod_{i\in J}a_i\Bigr)x .
\]
Let \(\mathcal J=\mathcal J_0\cup\{J\cup\{n+1\}:J\in\mathcal J_1\}\), writing
\(a_{n+1}\) for \(x\).  The two parts are disjoint, because members of the
second contain \(n+1\) and members of the first do not, and
\(J\mapsto J\cup\{n+1\}\) is injective; so the sum over \(\mathcal J\) is the
displayed sum and \(f=\Sigma_{\mathcal J}\).

If every member of \(\mathcal J\) has at most one element, then
\(f(\bar a)=c_0\oplus\bigoplus_ic_ia_i\) with \(c_0=1\) iff
\(\varnothing\in\mathcal J\) and \(c_i=1\) iff \(\{i\}\in\mathcal J\); so
\(f\in\mathsf L\).  The last sentence of the lemma is the contrapositive.
\end{proof}
\end{annexitem}

\begin{annexitem}{annex:p02-w02-p02-008-post}{thm:p02-w02-p02-008-criterion}{Post's criterion}
\begin{proof}
\emph{Part 1, necessity.}  Suppose \(C\subseteq K\) for one of the five
classes.  \(K\) is a proper clone (\cref{thm:p02-w02-p02-008-five-clones}),
so \(\langle C\rangle\subseteq K\neq\mathcal O\)
(\cref{thm:p02-w02-p02-008-generated}(2)).

\emph{Part 1, sufficiency.}  Choose \(f_0\in C\setminus\mathsf T_0\),
\(f_1\in C\setminus\mathsf T_1\), \(f_S\in C\setminus\mathsf S\),
\(f_M\in C\setminus\mathsf M\) and \(f_L\in C\setminus\mathsf L\), and let
\(K=\langle C\rangle\), a clone including \(C\).  There are exactly four
unary operations, since each is fixed by its two values: \(\kappa^1_0\),
\(\kappa^1_1\), \(e^1_1\) and \(\neg\).

\emph{Step 1: diagonals.}  For \(f\in K\cap\mathcal O_m\) the diagonal
\(f^\Delta=f\circ(e^1_1,\ldots,e^1_1)\in K\cap\mathcal O_1\) has
\(f^\Delta(x)=f(x,\ldots,x)\).  \(f_0^\Delta(0)=f_0(\bar0)=1\), so
\(f_0^\Delta\) is \(\kappa^1_1\) or \(\neg\); \(f_1^\Delta(1)=f_1(\bar1)=0\),
so \(f_1^\Delta\) is \(\kappa^1_0\) or \(\neg\).

\emph{Step 2: both unary constants lie in \(K\).}  If neither diagonal is
\(\neg\), they are \(\kappa^1_1\) and \(\kappa^1_0\), both in \(K\).
Otherwise \(\neg\in K\).  Since \(f_S\notin\mathsf S\) and \(f_S\) has some
arity \(n\), there is \(\bar a\in\mathbb B^n\) with
\(f_S(\bar a^{\,c})\neq1-f_S(\bar a)\), that is,
\(f_S(\bar a^{\,c})=f_S(\bar a)\).  For \(i\leq n\) let \(t_i=e^1_1\) if
\(a_i=1\) and \(t_i=\neg\) if \(a_i=0\); each \(t_i\in K\).  Let
\(u=f_S\circ(t_1,\ldots,t_n)\in K\).  Then \(t_i(1)=a_i\) and
\(t_i(0)=1-a_i\) for every \(i\), so \(u(1)=f_S(\bar a)\) and
\(u(0)=f_S(\bar a^{\,c})=f_S(\bar a)\).  So \(u=\kappa^1_b\) with
\(b=f_S(\bar a)\), and \(\neg\circ(u)=\kappa^1_{1-b}\in K\).

\emph{Step 3: negation lies in \(K\).}  By
\cref{thm:p02-w02-p02-008-tests}(2), since \(f_M\notin\mathsf M\) and
\(f_M\) has some arity \(n\), there are \(\bar a\) and \(i\) with
\(a_i=0\), \(f_M(\bar a)=1\) and \(f_M(\bar a[i{:=}1])=0\).  Let
\(t_i=e^1_1\) and \(t_j=\kappa^1_{a_j}\) for \(j\neq i\), all in \(K\) by
step~2, and \(w=f_M\circ(t_1,\ldots,t_n)\in K\).  Then
\(w(x)=f_M(\bar a[i{:=}x])\); as \(a_i=0\), \(w(0)=f_M(\bar a)=1\) and
\(w(1)=0\).  So \(w=\neg\in K\).

\emph{Step 4: conjunction lies in \(K\).}  The binary operations
\(\kappa^2_b=\kappa^1_b\circ(e^2_1)\) and \(\neg\circ(e^2_k)\)
(\(k=1,2\)) lie in \(K\).  Let \(n\) be the arity of \(f_L\).  By
\cref{thm:p02-w02-p02-008-anf}, \(f_L=\Sigma_{\mathcal J}\) for a set
\(\mathcal J\) of subsets of \(\{1,\ldots,n\}\), and since
\(f_L\notin\mathsf L\), \(\mathcal J\) has members with at least two
elements.  Among these choose \(I\) with the fewest elements, and choose
\(i\neq j\) in \(I\).  Let \(s_i=e^2_1\), \(s_j=e^2_2\),
\(s_k=\kappa^2_1\) for \(k\in I\setminus\{i,j\}\) and \(s_k=\kappa^2_0\)
for \(k\notin I\), and \(p=f_L\circ(s_1,\ldots,s_n)\in K\cap\mathcal O_2\).
Then \(p(x,y)=f_L(\bar a)\), where \(a_i=x\), \(a_j=y\), \(a_k=1\) for
\(k\in I\setminus\{i,j\}\) and \(a_k=0\) for \(k\notin I\).  Evaluate each
monomial \(\prod_{k\in J}a_k\), \(J\in\mathcal J\), at this \(\bar a\).
\begin{itemize}
\item If \(J\not\subseteq I\), \(J\) has a member \(k\notin I\) with
  \(a_k=0\), and the monomial is \(0\).
\item If \(J=I\), the monomial is \(xy\), since every other factor is \(1\).
\item If \(J\subsetneq I\), then \(J\) has fewer elements than \(I\), so by
  the choice of \(I\) it has at most one element.  The monomial is \(1\) if
  \(J=\varnothing\) or \(J=\{k\}\) with \(k\in I\setminus\{i,j\}\); it is
  \(x\) if \(J=\{i\}\) and \(y\) if \(J=\{j\}\).
\end{itemize}
Hence \(p(x,y)=xy\oplus\alpha x\oplus\beta y\oplus\gamma\), where
\(\alpha=1\) iff \(\{i\}\in\mathcal J\), \(\beta=1\) iff
\(\{j\}\in\mathcal J\), and \(\gamma\) is the parity of the number of members
of \(\mathcal J\) that are \(\varnothing\) or a singleton \(\{k\}\) with
\(k\in I\setminus\{i,j\}\).  Now
\[
\begin{aligned}
 p(x\oplus\beta,\,y\oplus\alpha)
 &=(x\oplus\beta)(y\oplus\alpha)\oplus\alpha(x\oplus\beta)\oplus
   \beta(y\oplus\alpha)\oplus\gamma\\
 &=(xy\oplus\alpha x\oplus\beta y\oplus\alpha\beta)\oplus
   (\alpha x\oplus\alpha\beta)\oplus(\beta y\oplus\alpha\beta)\oplus\gamma\\
 &=xy\oplus\alpha\beta\oplus\gamma,
\end{aligned}
\]
using \(z\oplus z=0\) to cancel \(\alpha x\), \(\beta y\) and two of the three
\(\alpha\beta\).  So \(xy=p(x\oplus\beta,y\oplus\alpha)\oplus\delta\) with
\(\delta=\alpha\beta\oplus\gamma\).  The arguments are \(e^2_1\) or
\(\neg\circ(e^2_1)\) according as \(\beta=0\) or \(1\), and \(e^2_2\) or
\(\neg\circ(e^2_2)\) according as \(\alpha=0\) or \(1\); adding
\(\delta=1\) is composing with \(\neg\).  So \(\land\in K\).

\emph{Step 5.}  \(\neg,\land\in K\), so
\(\langle\{\neg,\land\}\rangle\subseteq K\)
(\cref{thm:p02-w02-p02-008-generated}(3)), and
\(\langle\{\neg,\land\}\rangle=\mathcal O\)
(\cref{thm:p02-w02-p02-008-neg-and}).  So \(K=\mathcal O\).

\emph{Part 2.}  By \cref{thm:p02-w02-p02-008-formulas-clone}, \(C\) is
positive-arity functionally complete iff \(\langle C^+\rangle=\mathcal O\),
which by part~1 holds iff \(C^+\) meets the five conditions.  By
\cref{thm:p02-w02-p02-008-nullary}(2), \(C\) is nullary-inclusive
functionally complete iff it is positive-arity functionally complete and has
a nullary member.
\end{proof}
\end{annexitem}

\begin{annexitem}{annex:p02-w02-p02-008-maximal}{thm:p02-w02-p02-008-maximal}{The maximal clones}
\begin{proof}
\emph{1.}  Let \(K\) be a proper clone.  By
\cref{thm:p02-w02-p02-008-generated}(4), \(\langle K\rangle=K\neq\mathcal
O\), so by \cref{thm:p02-w02-p02-008-criterion}(1), \(K\subseteq\mathsf X\)
for at least one of the five classes \(\mathsf X\).

\emph{2.}  Let \(\mathsf X\) be one of the five classes; it is a proper clone
by \cref{thm:p02-w02-p02-008-five-clones}.  Let \(K\) be a clone with
\(\mathsf X\subsetneq K\).  Some \(f\in K\) lies outside \(\mathsf X\), so
\(K\not\subseteq\mathsf X\).  For each of the other four classes
\(\mathsf Y\), \(\mathsf X\not\subseteq\mathsf Y\)
(\cref{thm:p02-w02-p02-008-five-clones}), so some member of
\(\mathsf X\subseteq K\) lies outside \(\mathsf Y\) and
\(K\not\subseteq\mathsf Y\).  By \cref{thm:p02-w02-p02-008-criterion}(1),
\(\langle K\rangle=\mathcal O\), and \(\langle K\rangle=K\).  So the only
clone properly including \(\mathsf X\) is \(\mathcal O\), and \(\mathsf X\)
is maximal.

\emph{3.}  Each of the five is maximal by part~2.  Conversely, let \(K\) be
maximal.  It is proper, so \(K\subseteq\mathsf X\) for one of the five by
part~1.  If \(K\neq\mathsf X\), then \(\mathsf X\) is a clone properly
including \(K\) and different from \(\mathcal O\), against maximality; so
\(K=\mathsf X\).  The five classes are pairwise distinct and incomparable by
\cref{thm:p02-w02-p02-008-five-clones}, so there are exactly five maximal
clones.
\end{proof}
\end{annexitem}

\begin{annexitem}{annex:p02-w02-p02-008-finite}{thm:p02-w02-p02-008-finite}{A complete set has a complete subset of at most five members}
\begin{proof}
By \cref{thm:p02-w02-p02-008-criterion}(1), for each of the five classes
\(\mathsf X\) there is \(f_{\mathsf X}\in C\setminus\mathsf X\).  Let
\(D=\{f_{\mathsf T_0},f_{\mathsf T_1},f_{\mathsf S},f_{\mathsf M},
f_{\mathsf L}\}\subseteq C\), which has at most five members.  For each class
\(\mathsf X\), \(f_{\mathsf X}\in D\setminus\mathsf X\), so the criterion
gives \(\langle D\rangle=\mathcal O\).
\end{proof}
\end{annexitem}

\begin{annexitem}{annex:p02-w02-p02-008-nand-nor}{thm:p02-w02-p02-008-nand-nor}{NAND and NOR}
\begin{proof}
\emph{1.}  Each identity is checked from \(\mathrm{NAND}(x,y)=1\oplus xy\),
\(\mathrm{NOR}(x,y)=(1\oplus x)(1\oplus y)\), \(xx=x\) and \(x\oplus x=0\):
\(\mathrm{NAND}(x,x)=1\oplus x\);
\(\mathrm{NAND}(z,z)=1\oplus z\) with \(z=1\oplus xy\) gives \(xy\);
\(\mathrm{NAND}(x,1\oplus x)=1\oplus(x\oplus x)=1\); and
\(\mathrm{NAND}(1,1)=0\).  Further, \(\mathrm{NOR}(x,x)=1\oplus x\);
\(\mathrm{NOR}(1\oplus x,1\oplus y)=(1\oplus1\oplus x)(1\oplus1\oplus y)=xy\);
\(\mathrm{NOR}(x,1\oplus x)=(1\oplus x)x=x\oplus x=0\); and
\(\mathrm{NOR}(0,0)=1\).  In clone terms the identities say, for example,
\(\neg=\mathrm{NAND}\circ(e^1_1,e^1_1)\) and
\(\land=\mathrm{NAND}\circ(h,h)\) with
\(h=\mathrm{NAND}\circ(e^2_1,e^2_2)=\mathrm{NAND}\).

\emph{2.}  By part~1, \(\neg\) and \(\land\) are compositions of
\(\mathrm{NAND}\) with projections and with \(\mathrm{NAND}\) itself, so
they lie in the clone \(\langle\{\mathrm{NAND}\}\rangle\).  By
\cref{thm:p02-w02-p02-008-generated}(3) and
\cref{thm:p02-w02-p02-008-neg-and},
\(\mathcal O=\langle\{\neg,\land\}\rangle\subseteq
\langle\{\mathrm{NAND}\}\rangle\).  The same argument applies to
\(\mathrm{NOR}\): first \(\neg=\mathrm{NOR}\circ(e^1_1,e^1_1)\) lies in
\(\langle\{\mathrm{NOR}\}\rangle\), and then so does
\(\land=\mathrm{NOR}\circ(\neg\circ(e^2_1),\neg\circ(e^2_2))\).  The operation set of the connective
set \(\{\mathrm{NAND}\}\) is \(\{\mathrm{NAND}\}\), so by
\cref{thm:p02-w02-p02-008-formulas-clone} the connective set is
positive-arity functionally complete; likewise for \(\{\mathrm{NOR}\}\).

\emph{3.}  Neither set has a nullary member, so
\cref{thm:p02-w02-p02-008-nullary}(1) applies.  The operation set of
\(\{\mathrm{NAND},\bottom\}\) is \(\{\mathrm{NAND},\kappa^1_0\}\), which
includes \(\{\mathrm{NAND}\}\), so it generates \(\mathcal O\) by part~2 and
\cref{thm:p02-w02-p02-008-generated}(3); \(\bottom\) is nullary, so
\cref{thm:p02-w02-p02-008-nullary}(2) gives nullary-inclusive completeness.
The closed formulas representing the two constants are \(\bottom\) and
\(\mathrm{NAND}(\bottom,\bottom)\), whose value is \(\mathrm{NAND}(0,0)=1\).
For \(\{\mathrm{NOR},\bottom\}\) the argument is the same, with the closed
formulas \(\bottom\) and \(\mathrm{NOR}(\bottom,\bottom)\), of value
\(\mathrm{NOR}(0,0)=1\).
\end{proof}
\end{annexitem}

\begin{annexitem}{annex:p02-w02-p02-008-sheffer}{thm:p02-w02-p02-008-sheffer}{The singly complete operations of arity at most two}
\begin{proof}
\emph{Unary.}  The four unary operations each lie in one of the classes, by
Table~\ref{tab:p02-w02-p02-008-profile}: \(\kappa^1_0,e^1_1\in\mathsf T_0\),
\(\kappa^1_1\in\mathsf T_1\), \(\neg\in\mathsf S\).  So
\(\langle\{u\}\rangle\neq\mathcal O\) by
\cref{thm:p02-w02-p02-008-criterion}(1).

\emph{Binary.}  Let \(f\in\mathcal O_2\) with
\(\langle\{f\}\rangle=\mathcal O\).  By the criterion \(f\notin\mathsf T_0\)
and \(f\notin\mathsf T_1\), so \(f(0,0)=1\) and \(f(1,1)=0\).  If
\(f(0,1)=1\) and \(f(1,0)=0\), then \(f(x,y)=1-x\) at all four inputs, and
\(f(\bar a^{\,c})=1-(1-a_1)=1-f(\bar a)\), so \(f\in\mathsf S\), which is
impossible.  If \(f(0,1)=0\) and \(f(1,0)=1\), then \(f(x,y)=1-y\), which is
self-dual in the same way.  So \(f(0,1)=f(1,0)\): if both are \(1\), \(f\)
has values \(1,1,1,0\) and is \(\mathrm{NAND}\); if both are \(0\), \(f\) has
values \(1,0,0,0\) and is \(\mathrm{NOR}\).  Conversely both generate
\(\mathcal O\) by \cref{thm:p02-w02-p02-008-nand-nor}(2).
\end{proof}
\end{annexitem}

\begin{annexitem}{annex:p02-w02-p02-008-bases}{thm:p02-w02-p02-008-bases}{Negative and positive bases}
\begin{proof}
For a set \(C\) of operations and a class \(\mathsf X\) among the five,
\(C\subseteq\mathsf X\) iff every member of \(C\) lies in \(\mathsf X\), and
by \cref{thm:p02-w02-p02-008-generated}(5) and
\cref{thm:p02-w02-p02-008-five-clones} this holds iff
\(\langle C\rangle\subseteq\mathsf X\).  So the classes including \(C\) are
the intersection of the rows of its members in
Table~\ref{tab:p02-w02-p02-008-profile}.  If \(\langle C\rangle\subseteq
\mathsf X\) and \(g\notin\mathsf X\), then \(g\notin\langle C\rangle\); for a
connective set whose operation set is \(C\), no formula represents \(g\), by
\cref{thm:p02-w02-p02-008-formulas-clone}.  Each item follows.
\begin{enumerate}
\item \(\land\) and \(\lor\) both lie in exactly
  \(\mathsf T_0,\mathsf T_1,\mathsf M\).  \(\neg\) and \(\kappa^1_1\) lie
  outside \(\mathsf T_0\), and \(\kappa^1_0\) outside \(\mathsf T_1\).
\item \(\land,\lor\) give \(\mathsf T_0,\mathsf T_1,\mathsf M\);
  \(\kappa^1_0\) removes \(\mathsf T_1\) and \(\kappa^1_1\) removes
  \(\mathsf T_0\), leaving \(\mathsf M\), to which all four belong.  \(\neg\)
  lies outside \(\mathsf M\).  The operation set of
  \(\{\land,\lor,\bottom,\topv\}\) is this \(C\).
\item \(\neg\) lies in exactly \(\mathsf S,\mathsf L\); \(\leftrightarrow\) in
  exactly \(\mathsf T_1,\mathsf L\); \(\oplus\) in exactly
  \(\mathsf T_0,\mathsf L\).  Both intersections are \(\{\mathsf L\}\).  As
  \(1\oplus x\oplus y=1\oplus(x\oplus y)\),
  \(\leftrightarrow=\neg\circ(\oplus)\) and \(\oplus=\neg\circ(\leftrightarrow)\),
  so each set is included in the clone generated by the other, and
  \cref{thm:p02-w02-p02-008-generated}(3) gives equality of the two clones.
  \(\land\), \(\lor\) and \(\to\) lie outside \(\mathsf L\).
\item \(\to\) lies in exactly \(\mathsf T_1\); \(\neg\) and \(\kappa^1_0\)
  lie outside it.
\item \(\land\) lies in exactly \(\mathsf T_0,\mathsf T_1,\mathsf M\) and
  \(\oplus\) in exactly \(\mathsf T_0,\mathsf L\), with intersection
  \(\{\mathsf T_0\}\).  \(\neg\), \(\to\), \(\leftrightarrow\) and
  \(\kappa^1_1\) lie outside \(\mathsf T_0\).
\item \(\neg\) lies in exactly \(\mathsf S,\mathsf L\) and
  \(\operatorname{maj}\) in exactly
  \(\mathsf T_0,\mathsf T_1,\mathsf S,\mathsf M\), with intersection
  \(\{\mathsf S\}\).  \(\land\), \(\lor\), \(\kappa^1_0\) and \(\kappa^1_1\)
  lie outside \(\mathsf S\).
\item \(\to\) lies only in \(\mathsf T_1\), and \(\kappa^1_0\) is not in
  \(\mathsf T_1\); so no class includes \(\{\to,\kappa^1_0\}\).  \(\land\)
  and \(\oplus\) share only \(\mathsf T_0\), and \(\kappa^1_1\) is not in
  \(\mathsf T_0\); so no class includes \(\{\land,\oplus,\kappa^1_1\}\).
  These are the operation sets of \(\{\to,\bottom\}\) and
  \(\{\land,\oplus,\topv\}\), each of which has a nullary member, so
  \cref{thm:p02-w02-p02-008-criterion}(2) gives completeness in both senses.
\end{enumerate}
\end{proof}
\end{annexitem}

\begin{annexitem}{annex:p02-w02-p02-008-alt}{thm:p02-w02-p02-008-alternating}{Alternating means affine}
\begin{proof}
Write \(\bar a^{(i)}=\bar a[i{:=}1-a_i]\).

Let \(f(\bar a)=c_0\oplus\bigoplus_kc_ka_k\).  Replacing \(a_i\) by
\(1-a_i=1\oplus a_i\) changes the term \(c_ia_i\) into
\(c_i\oplus c_ia_i\), so \(f(\bar a^{(i)})=f(\bar a)\oplus c_i\) for every
\(\bar a\).  If \(c_i=0\), the first alternative holds; if \(c_i=1\), the
second.

Conversely, suppose that for each \(i\) one of the alternatives holds.  They
cannot both hold, because \(\mathbb B^n\) is nonempty and \(f(\bar a)\neq
1-f(\bar a)\).  Let \(d_i=0\) if the first holds and \(d_i=1\) if the second
does; then \(f(\bar a^{(i)})=f(\bar a)\oplus d_i\) for every \(\bar a\).  We
show \(f(\bar a)=f(\bar0)\oplus\bigoplus_kd_ka_k\) by induction on the number
of coordinates of \(\bar a\) equal to \(1\).  For \(\bar0\) the sum is
\(f(\bar0)\).  Otherwise pick \(i\) with \(a_i=1\) and let
\(\bar b=\bar a[i{:=}0]\), which has one fewer coordinate equal to \(1\) and
satisfies \(\bar b^{(i)}=\bar a\).  By the induction hypothesis and the
displayed property,
\(f(\bar a)=f(\bar b)\oplus d_i=f(\bar0)\oplus\bigoplus_kd_kb_k\oplus d_i\),
and since \(b_k=a_k\) for \(k\neq i\), \(b_i=0\) and \(a_i=1\), the right side
is \(f(\bar0)\oplus\bigoplus_kd_ka_k\).  So \(f\in\mathsf L\).
\end{proof}
\end{annexitem}
